% ============================================================
% ch25.tex —— 第 25 章 曲线的性格：曲率
% ============================================================
\chapter{曲线的性格：曲率}

篇五换装。前三篇我们研究"数"与"变化"，现在把镜头对准\textbf{形状}——第一站：曲线。导数（篇三）已经量过曲线的\textbf{倾斜}；这一章发明新工具量它的\textbf{弯曲}。这个工具将一路升级，最终让蚂蚁量出宇宙的形状（第 27、28 章）。

\begin{kaiwei}
高速公路的警示牌写着"前方急转弯，限速 40"。"急"是司机的体感——能把它变成一个\textbf{数}吗？急到什么程度算急，急多少，说出来要带单位。
\end{kaiwei}

\section{先给"弯"定标尺}

任何测量先要定\textbf{零点}和\textbf{单位}。

\textbf{零点：直线}。直路不弯，规定直线的"弯度"为 $0$——没有争议。

\textbf{标准件：圆}。圆是"弯得最匀"的曲线——每一点、每个方向的弯法完全相同（转一圈回到原处，处处一样，均匀性无出其右）。用圆当"弯度标准砝码"，直觉立刻给出方向：\textbf{半径越小，弯得越急}。游乐场环形小跑道（半径 $5$ 米）弯得人头晕，高速大圆弧（半径 $2000$ 米）你几乎无感——同样的"转一圈"，小圆挤在几米里完成，大圆摊开几公里。

"急"与半径\textbf{反向}相关，最自然的标尺就是取倒数：

\begin{dingyi}[圆的弯度（曲率的定义）]
半径 $r$ 的圆，弯度规定为
\[
\kappa \;=\; \frac{1}{r} \qquad (\kappa \text{ 读 kappa}) .
\]
单位：$1/$米。
\end{dingyi}

这个定义不是拍脑袋，它精确捕捉了司机的体感——\textbf{单位路程上的转头量}。看推导：开车沿半径 $r$ 的圆兜一圈，方向盘累计转过的角度是 $360^\circ = 2\pi$ 弧度，跑过的路程是周长 $2\pi r$。单位路程分到的转头量：
\[
\frac{2\pi \text{ 弧度}}{2\pi r \text{ 米}} = \frac{1}{r} \text{ 弧度/米} = \kappa .
\]
（弧度的定义在此严丝合缝：单位圆上走 $1$ 米弧长，转角 $1$ 弧度。）所以"半径 $50$ 米的匝道，$\kappa = 0.02$ /米"翻译成人话：\textbf{每前进 $1$ 米，车头方向扳回 $0.02$ 弧度}（约 $1.15^\circ$）。半径 $2000$ 米的主线，$\kappa = 0.0005$/米——差 $40$ 倍，体感与数字对上了。

\section{一般曲线：每一点都有自己的弯度}

圆处处一样弯，但抛物线不是：顶点处最陡峭地转弯，往两边越来越"直"（最后几乎躺平）。给一般曲线定义弯度，用微积分的标准姿势——\textbf{局部化}：在每个点，找一个\textbf{最贴合的圆}（土名：贴身圆），看它与曲线在该点"方向相同、弯曲趋势相同"——然后用贴身圆的弯度作为曲线在该点的弯度：
\[
\text{曲线在一点的曲率} \;\kappa\; =\; \frac{1}{r_{\text{贴身}}} .
\]

\begin{faming}
\textbf{土名}：弯度、转弯强度 $\to$ \textbf{正式名}：\textbf{曲率}（curvature），$\kappa$。贴身圆的正式名是\textbf{密切圆}（osculating circle，拉丁语 osculari"亲吻"——它在该点与曲线贴得最紧，"吻"到二阶）；密切圆的圆心与半径分别叫\textbf{曲率中心}与\textbf{曲率半径}（$= 1/\kappa$）。曲率是微分几何的第一块砖：第 26 章升到曲面，第 27 章让蚂蚁测它，第 28 章用它给宇宙分类。
\end{faming}

\section{曲率用导数算}

贴身圆怎么找？不必真画圆——微积分直接给公式（对 $y = f(x)$）：
\[
\boxed{\;\kappa = \frac{|y''|}{\left(1 + y'^2\right)^{3/2}}\;}
\]
分子 $y''$ 是\textbf{导数的导数}——倾斜的变化快慢（篇三的伏笔兑现：一阶导管斜，二阶导管转弯）；分母是斜率修正项：路越斜，同样的 $y''$ 分摊在更长的弧上，每米分到的转头就少（除以"弧长膨胀因子" $(1+y'^2)^{3/2}$——勾股定理斜边因子的二分之三次方）。

实战两发，各带完整计算：

\paragraph{发一：抛物线 $y = x^2$ 的顶点。}$y' = 2x$，顶点 $x = 0$ 处 $y' = 0$，$y'' = 2$。代入：
\[
\kappa(0) = \frac{|2|}{(1+0)^{3/2}} = 2 .
\]
顶点曲率 $2$（坐标单位取米时，即 $2$/米；贴身圆半径 $\tfrac12$ 米——抛物线顶点比半径半米的圆还急）。再算个旁点对比：$x = 1$ 处 $y' = 2$，$y'' = 2$：
\[
\kappa(1) = \frac{2}{(1+4)^{3/2}} = \frac{2}{5\sqrt5} \approx 0.179 ,
\]
离顶点一步之遥，曲率掉到不足十分之一——\textbf{抛物线顶点最弯、两侧迅速躺平}，图像脑补即可验证。

\paragraph{发二：圆 $x^2 + y^2 = r^2$ 处处 $\kappa = \tfrac1r$。}取上半圆 $y = \sqrt{r^2 - x^2}$。求导（链式法则，内层导 $-2x$）：
\[
y' = \frac{-x}{\sqrt{r^2-x^2}}, \qquad
y'' = -\frac{r^2}{\left(r^2 - x^2\right)^{3/2}}
\quad\text{(}\sqrt{r^2 - x^2}\text{ 约掉后正好剩下 }r^2\text{，请亲手通分验算)} .
\]
代入公式，分子分母的根号精确相消（算一算就看到 $r^2$ 自己冒出来），得 $\kappa = \tfrac1r$，\textbf{与 $x$ 无关}——处处相同。\textbf{圆的曲率是常数}，这正是当初敢拿它当"标准砝码"的资格证：公式自动归还了定义，机器自洽。

\section{为什么司机和过山车工程师在乎}

\paragraph{物理兑现：向心力。}质量 $m$、速度 $v$ 的物体走曲率 $\kappa$ 的弯，需要向心力
\[
F = m\,v^2\kappa \qquad (\text{即熟悉的 } \tfrac{mv^2}{r}).
\]
代入数字感受量级：一辆 $1.5$ 吨轿车以 $120$ km/h（$33.3$ m/s）过半径 $50$ 米的匝道（$\kappa = 0.02$）：
\[
F = 1500 \times 33.3^2 \times 0.02 \approx 33\,000 \text{ 牛} \approx 3.4 \text{ 吨力}.
\]
而轮胎能给的极限摩擦力：车重 $1500\times9.8 = 14\,700$ 牛，干燥路面摩擦系数约 $0.8$，上限 $0.8\times14\,700 \approx 11\,800$ 牛——\textbf{远小于所需的 $33\,000$ 牛，车必滑出弯道}。所以匝道限速 $40$ km/h（$11.1$ m/s）：$F = 1500\times11.1^2\times0.02 \approx 3\,700$ 牛，仅用掉摩擦上限的三成，安全。雨后摩擦减半，限速还得再砍——"急弯限速"牌上的每一个数字，背后都是这条公式。

\paragraph{工程兑现：曲率连续。}过山车轨道的拼接有一条铁律：\textbf{连接处曲率必须连续}（土名：转弯不跳档）。若直线段（$\kappa = 0$）直接接圆弧（$\kappa = 1/r \ne 0$），曲率瞬间跳变，向心力瞬间从 $0$ 跳到 $mv^2/r$——乘客的颈椎第一个抗议（工程术语"冲击"）。解法：直线与圆弧之间插入\textbf{缓和曲线}——一段曲率从 $0$ \textbf{线性缓升}到 $1/r$ 的过渡曲线（铁路上的"回旋曲线"/放射螺线，曲率与弧长成正比）。高铁进出弯道时的轻微"渐倾感"、过山车起始爬升的丝滑，都是设计师在给你的人生做\textbf{曲率的积分平滑}。

\begin{cidian}
土名对照：弯度 $=$ \textbf{曲率} $\kappa$；贴身圆 $=$ \textbf{密切圆}；转弯不跳档 $=$ \textbf{曲率连续性}（$G^2$ 连续）；缓和段 $=$ \textbf{回旋曲线}（clothoid）。物理链条一页看全：$y'$ 给斜率，$y''$（配合修正）给曲率，$v^2\kappa$ 给向心加速度。
\end{cidian}

\begin{dongshou}
\begin{itemize}
  \yisi 心算下列曲线顶点处的曲率（顶点处 $y' = 0$，分母为 $1$，$\kappa = |y''|$）：$y = x^2$（答 $2$）；$y = 3x^2$（答 $6$——图像被纵向拉伸三倍，弯度也三倍）；$y = x^2 + 100x + 7$（答 $2$——\textbf{平移不改变曲率}：配方可得顶点挪到别处，弯度原地不动。几何量不认坐标的脸色，这是第 27 章"内在量"的伏笔）。
  \yier 完整计算圆 $r = 5$ 的曲率公式验证：$y = \sqrt{25 - x^2}$，求 $y'$、$y''$（答 $y'' = -25/(25-x^2)^{3/2}$），代入曲率公式并化简（分子分母相消，答 $\tfrac15$，与 $x$ 无关）。
  \yier 骑车绕半径 $5$ 米的小转盘，速度 $3$ m/s：向心加速度 $v^2\kappa = 1.8$ m/s$^2$，不到重力的两成，安全。速度提到 $10$ m/s：$v^2\kappa = 20$ m/s$^2 = 2g$——摔。复算这两组数，并求"临界速度"（向心加速度恰为 $g\times$摩擦系数 $0.6$ 时，$v = \sqrt{0.6\times9.8\times5} \approx 5.4$ m/s $\approx 19.5$ km/h）。
\end{itemize}
\end{dongshou}

\begin{shaonao}
\textbf{"曲率公式为什么长这样"的推导骨架}（跟完这几步，公式就不再是天书）：

(1) 曲线的"方向角" $\theta = \arctan(y')$——切线与横轴的夹角。曲率的本义是\textbf{方向角对弧长的变化率}：$\kappa = \left|\dfrac{d\theta}{ds}\right|$。

(2) 两件微积分存货：$ds = \sqrt{1 + y'^2}\,dx$（弧长微分——一小段斜路的水平投影 $dx$、竖直增量 $y'dx$，勾股定理得斜边）；$\theta = \arctan(y')$ 对 $x$ 求导（链式）：$\dfrac{d\theta}{dx} = \dfrac{y''}{1+y'^2}$。

(3) 除法衔接：$\kappa = \left|\dfrac{d\theta/dx}{ds/dx}\right| = \dfrac{|y''|/(1+y'^2)}{\sqrt{1+y'^2}} = \dfrac{|y''|}{(1+y'^2)^{3/2}}$。

请在 (2)(3) 之间各补一行计算，体会整个公式只是\textbf{"转头速度 ÷ 前进速度"}换了个记号。再检验一发：对圆 $\theta$ 转一圈 ($2\pi$) 弧长 $2\pi r$，$\kappa = 2\pi/2\pi r = 1/r$——定义、公式、直觉三方对齐。\textbf{能自己推一遍的公式，才是自己的公式}。
\end{shaonao}

\begin{chengshi}
曲率公式的完整推导（含参数方程与空间曲线版本、密切圆的存在唯一性）见 do Carmo《曲线与曲面的微分几何》第一章；$|y'| \ll 1$ 时分母 $\approx 1$、$\kappa \approx |y''|$（"小斜率近似"），顶点计算我们一直这么用。向心力公式 $F = mv^2/r$ 本书直接引用（高中内容，严格推导需向量微积分）。三维空间曲线的曲率同款定义（方向角对弧长的变化率），但空间多了一个"往哪边扭"的量（挠率）——那是大学的故事。
\end{chengshi}

曲线的弯度量完了。升维：曲面。而曲面上的"直线"是什么？曲线的"直"有导数为零作证，曲面的"直"谁作证？下一章，去地球仪上找答案——顺便眼睁睁看着初中几何的最大公理崩塌。
